\section{The Fathers: created light and received holiness}
\label{sec:fathers}

The angels belong to the city whose happiness is communion with God. Augustine begins his account of their origin from that membership: the holy angels are the part of the heavenly city that has never gone into exile. Their superiority to us is therefore neither independence from God nor a separate destiny. They possess already the blessedness toward which the pilgrim Church travels. This places angelology inside the history of salvation, where created excellence and dependence upon grace belong together.\footnote{Augustine, \emph{City of God} XI.9, trans. Marcus Dods, vol.~I (Edinburgh: T.~\& T.~Clark, 1871).}

\subsection{Augustine: the light that receives light}

Genesis does not expressly name the creation of angels. Augustine acknowledges this before proposing that the light of the first day signifies their creation and illumination. The certainty that they are God's creatures comes also from the praise of the Creator in Psalm 148; the identification of the first light with the angels is his interpretation of Genesis. In XI.9, the angels become light through participation in the eternal Word. Those who turn away become darkness through their own deprivation. Evil is thus a loss in a good creature, not a second created substance or an independent origin of being.


Augustine's argument begins from worship. The angels stand among the works summoned to praise their maker in the canticle of the three children and in Psalm 148. Their praise is itself testimony to their created condition: they bless the one whose command gave them being. When he connects them with the first light, Augustine reads creation through this already established relation between the giver and the gift. The light becomes a particularly apt image because its excellence is real and its dependence continuous. An angel shines by receiving the Word who made it. To turn away is to lose the rectitude of that participation, even while the nature that God created remains good.

The contrast between light and darkness accordingly distinguishes wills as well as states of knowledge. The faithful angels cling to their common good; the apostate spirits seek themselves apart from him. The spiritual creature does not need a body in order to be capable of pride. Its own natural beauty can become an occasion of disorder if it loves that beauty as an independent end. Augustine's account thus places the decisive division within the response of the creature, rather than between one kind of material made by God and another made by an enemy. Both the holy and the fallen spirits bear witness to the Creator: the former through the right use of his gifts, the latter through the goodness of the nature whose order their will has violated.\footnote{Augustine, \emph{City of God} XI.15--17.}

Augustine therefore distinguishes the scale of nature from the scale of justice. In XI.16 angels stand above mortal humanity by their natural excellence, yet a good human being is better than a wicked angel. The will's adherence to God determines the creature's moral good. Knowledge and power become blessed through their right order to the one who gives them.

In XI.10 Augustine explains why even an immortal angel remains unlike the divine simplicity. A creature has gifts it is not: wisdom, goodness, and blessedness belong to it by participation. God is not a subject subsequently furnished with these perfections. Angelic incorporeality therefore does not erase the distinction between having life and being the source of life. Aquinas develops that distinction when he denies that an angel's act of understanding is its substance or its existence (\ST{I q.~54 aa.~1--2}). The absence of bodily parts is not the absence of every created composition.


This distinction is also an account of angelic happiness. Augustine's examples of a vessel and its contents, air and its light, and a mind and its wisdom show what it means for a creature to receive a perfection. The blessed angel's wisdom is no passing illusion: God makes the creature truly wise. Yet the wisdom points beyond the subject that possesses it. In God, life and wisdom are not gifts added to an already existing subject; he is the living good from whom their created participation comes. The angel's happiness therefore consists in adhering to a good greater than itself. Its freedom is perfected in that adherence, and its praise confesses the truth of its own being.

Augustine's account of knowledge adds a second distinction. Angels know creatures in the eternal Word, and also in the creatures' own being. In XI.29 he compares knowing a figure in the art that produces it with knowing the figure drawn on paper. The creature's own reality is good, but knowledge of it in its divine cause is clearer. When the knowledge of created things returns to the Creator in praise, evening gives way to morning. Aquinas receives this interpretation in his treatment of morning and evening knowledge (\ST{I q.~58 aa.~6--7}). These terms describe modes of angelic knowing, not the alternation of sunlight and darkness in an angel's environment.


The angels' knowledge of themselves follows the same pattern. They know themselves more perfectly in the Word than through their own created being considered by itself. Their highest knowledge is thus inseparable from humility: they understand their excellence by understanding its source. Augustine includes their knowledge of the Trinity in this account. The Father, Son, and Holy Spirit are present to their understanding as the one God, and the eternal reasons of creatures are seen in the Wisdom through whom those creatures exist. The celestial city is united by this common truth, not merely by the coordinated labor of many powerful individuals. Its contemplation issues in praise, and its praise rightly orders what it knows.\footnote{Augustine, \emph{City of God} XI.9--10, 29.}

The same theological distinction governs perseverance. Augustine knows that the holy angels will not fall, but XI.13 considers alternative explanations of the relation between their original condition, the fall of others, and assurance of endless blessedness. His argument rejects an evil nature created by an evil source while leaving parts of the chronology under discussion. Aquinas's later determination about the first acts of angels must accordingly be read as a determination within this inherited discussion, not as a chronology already supplied in every detail by Augustine.\footnote{Augustine, \emph{City of God} XI.9--10, 13, 29; compare \ST{I q.~62 aa.~1, 8; q.~63 aa.~5--6}.}

\subsection{Basil: the Spirit perfects the heavenly hosts}

Basil considers the angels while defending the inseparability of the Holy Spirit from the Father and the Son. In \emph{On the Holy Spirit} XVI.37--38, their sanctification provides an argument for the Spirit's divine dignity. The angels are holy through the Spirit; the Spirit is not one more recipient within their company. Basil speaks of the Father's originating causality, the Son's creative operation, and the Spirit's perfecting presence, immediately excluding the suggestion of three independent first principles or a defective work completed by another agent. Their creation and holiness manifest the one divine action.\footnote{Basil of Caesarea, \emph{On the Holy Spirit} XVI.37--38, trans. Blomfield Jackson, \emph{NPNF}, second series, vol.~VIII (Edinburgh: T.~\& T.~Clark, 1895), CCEL transcription.}

Angelic perfection, for Basil, includes sanctification and abiding in it. Prophecy, the announcement of mysteries, the praise sung by seraphim, and the concord of the heavenly hosts all depend upon the Spirit. Gabriel's knowledge of the future is a received gift, not a natural possession of every possible event. This is a theology of sustained dependence: angelic order itself would dissolve without the divine giver of holiness. Aquinas's distinction between natural knowledge and revealed mysteries, and his teaching that an angel needs grace to attain supernatural beatitude, answer related questions with a more differentiated account of faculties and ends (\ST{I q.~57 aa.~3, 5; q.~62 a.~2}).


Basil's illustration is the iron made hot by fire. The heat really belongs to the iron, while iron and fire remain distinguishable. So the sanctification of the angels is truly theirs without being identical to the substance they received in creation. This provides the force of his argument against treating the Spirit as one of the ministering powers. The angels' excellence is a communicated holiness; the Spirit is its giver. The dependence appears at the summit of creation as surely as among the human recipients of spiritual gifts. Holiness does not cease to be received because its recipient is exalted.

Basil then draws the whole activity of heaven into the argument. Without the Spirit, angels could not give glory to God at the Nativity; the seraphim could not sing their threefold acclamation of holiness; the heavenly powers could not behold the Father's face and live their blessed life. Their concord resembles the order of an army under its commander and the harmony of a chorus under its leader. These comparisons disclose an actively sustained union. The Spirit supplies the light through which the mind sees and the holiness through which free spirits remain ordered to the good. Withdraw that gift in thought, and the supposed self-sufficiency of the heavenly world vanishes.

The resulting vision joins worship and service. A spirit's announcement of a mystery and its praise of the divine giver are two expressions of the same received illumination. Gabriel teaches because he has first received; the seraph burns with a holiness that comes from God; each order's distinct excellence contributes to the concord of the whole. The doctrine of angels thereby serves Basil's confession of the Trinity. Their splendor makes the Spirit's sanctifying power manifest, and their ordered freedom displays the goodness of the one divine action.\footnote{Basil, \emph{On the Holy Spirit} XVI.38.}

Basil also calls the angelic substance, tentatively, an aerial spirit or immaterial fire. His principal contrast here is between the created heavenly powers and the sanctifying Spirit: the powers receive holiness, while the Spirit sanctifies them. Aquinas determines the further question whether an angel contains any matter in \ST{I q.~50 a.~2}. Basil's account establishes the angels' dependence upon God in the excellence of their holiness and the order of their common life.
